\chapter{Estado del Arte}

Con el propósito de comprender el estado actual de la investigación relacionada con el uso de LLMs en el análisis de logs de seguridad, esta sección presenta la síntesis y análisis de los estudios seleccionados mediante el proceso de revisión sistemática descrito previamente. El análisis se estructura en función de las preguntas de investigación definidas, permitiendo examinar las principales técnicas, enfoques, capacidades y desafíos reportados en la literatura reciente. A partir de la evidencia recopilada, se identifican tendencias de desarrollo, fortalezas y limitaciones de las soluciones existentes, así como los vacíos de investigación que aún persisten en el área. Estos hallazgos constituyen la base para fundamentar la propuesta presentada en este seminario, orientada a abordar las brechas detectadas en el análisis semántico y la reconstrucción de incidentes de seguridad mediante arquitecturas basadas en modelos de lenguaje.

%%-------------------------------------------------------------
\section{Panorama bibliométrico del corpus}
%%-------------------------------------------------------------

Tras el análisis aplicado, respecto a la distribución geográfica, en la Figura~\ref{fig:publicaciones_por_pais} se puede apreciar una producción científica sobre inteligencia artificial y modelos de lenguaje aplicados al análisis de logs y la ciberseguridad, concentrada en países con ecosistemas tecnológicos consolidados. Estados Unidos aparece como el principal contribuyente, seguido de China. Asimismo, se observa una participación relevante de países europeos como Países Bajos, Suiza, Reino Unido y Alemania, junto con representantes de Asia y Oceanía. Esto evidencia una tendencia global hacia el desarrollo de esta área, aunque con una marcada concentración en un número reducido de países.


\begin{figure}[H]
    \centering
    \includegraphics[width=1\textwidth]{Resource/papersXpais_barras.png}
    \caption{Distribución geográfica de la producción científica en estudios
    sobre modelos de lenguaje de gran escala aplicados al análisis de logs y
    la ciberseguridad. Fuente: Elaboración propia.}
    \label{fig:publicaciones_por_pais}
\end{figure}

Por otro lado, la distribución temporal de las publicaciones presentada en la Figura~\ref{fig:publicaciones_por_año} muestra una clara tendencia de crecimiento entre 2023 y 2025. La cantidad de estudios aumenta desde 30 publicaciones en 2023 hasta 72 en 2024, alcanzando un máximo de 164 publicaciones en 2025. Este comportamiento evidencia un creciente interés de la comunidad científica por la aplicación de modelos de lenguaje de gran escala en el análisis de logs y la ciberseguridad. Aunque en 2026 se observa una disminución respecto al año anterior, con 91 publicaciones registradas, este valor continúa siendo considerablemente superior al de los primeros años analizados. En conjunto, estos resultados confirman que se trata de un área de investigación emergente que ha experimentado una expansión significativa en los últimos años.

\begin{figure}[H]
    \centering
    \includegraphics[width=1\textwidth]{Resource/paperxano.png}
    \caption{Distribución temporal de las publicaciones científicas sobre modelos de lenguaje de gran escala aplicados al análisis de logs y la ciberseguridad durante el período 2023–2026. Fuente: Elaboración propia.}
    \label{fig:publicaciones_por_año}
\end{figure}


%%-------------------------------------------------------------
\section{Hallazgos por categoría de análisis}
%%-------------------------------------------------------------
 
Esta sección presenta un análisis organizado según las categorías
de los trabajos del corpus. La Tabla~\ref{tab:hallazgos_parte1}
y la Tabla~\ref{tab:hallazgos_parte2} identifican los modelos y
técnicas empleados, los datasets utilizados, las características
de preprocesamiento e integración de conocimiento, y las
principales métricas de evaluación reportadas en la literatura.
 
% ============================================================
% Tabla 7.1 — Hallazgos parte 1 (papers 1–26)
% Columnas: Modelos/Técnicas | Datasets | Preprocesamiento/KB | Métricas
% ============================================================
\begin{longtable}{p{2.2cm}*{5}{>{\centering\arraybackslash\small}p{0.4cm}}*{5}{>{\centering\arraybackslash\small}p{0.4cm}}*{3}{>{\centering\arraybackslash\small}p{0.4cm}}*{4}{>{\centering\arraybackslash\small}p{0.6cm}}}

\caption{Hallazgos por categoría de análisis en artículos seleccionados
         (trabajos 1--26 de 52).}
\label{tab:hallazgos_parte1} \\
 
\toprule
& \multicolumn{5}{c}{\textbf{Modelos / Técnicas}} 
& \multicolumn{5}{c}{\textbf{Datasets}}
& \multicolumn{3}{c}{\textbf{Preproc./KB}}
& \multicolumn{4}{c}{\textbf{Métricas}} \\
 
\cmidrule(lr){2-6}\cmidrule(lr){7-11}\cmidrule(lr){12-14}\cmidrule(lr){15-18}
 
\textbf{ID} &
\rotatebox{90}{Fine-tuning} &
\rotatebox{90}{Emb. BERT} &
\rotatebox{90}{LLM/Prompting} &
\rotatebox{90}{ML Clásico} &
\rotatebox{90}{Agente/ReAct} &
\rotatebox{90}{HDFS} &
\rotatebox{90}{BGL} &
\rotatebox{90}{Thunderbird} &
\rotatebox{90}{UWF-Zeek} &
\rotatebox{90}{Logs Web/Red} &
\rotatebox{90}{Log Parsing} &
\rotatebox{90}{RAG / Base KB} &
\rotatebox{90}{MITRE ATT\&CK} &
\rotatebox{90}{Precision} &
\rotatebox{90}{Recall} &
\rotatebox{90}{F1-Score} &
\rotatebox{90}{Accuracy} \\
\midrule
\endfirsthead
 
\multicolumn{18}{c}{\tablename~\thetable{} (continuación)} \\
\toprule
& \multicolumn{5}{c}{\textbf{Modelos / Técnicas}} 
& \multicolumn{5}{c}{\textbf{Datasets}}
& \multicolumn{3}{c}{\textbf{Preproc./KB}}
& \multicolumn{4}{c}{\textbf{Métricas}} \\
\cmidrule(lr){2-6}\cmidrule(lr){7-11}\cmidrule(lr){12-14}\cmidrule(lr){15-18}
\textbf{ID} &
\rotatebox{90}{Fine-tuning} &
\rotatebox{90}{Emb. BERT} &
\rotatebox{90}{LLM/Prompting} &
\rotatebox{90}{ML Clásico} &
\rotatebox{90}{Agente/ReAct} &
\rotatebox{90}{HDFS} &
\rotatebox{90}{BGL} &
\rotatebox{90}{Thunderbird} &
\rotatebox{90}{UWF-Zeek} &
\rotatebox{90}{Logs Web/Red} &
\rotatebox{90}{Log Parsing} &
\rotatebox{90}{RAG / Base KB} &
\rotatebox{90}{MITRE ATT\&CK} &
\rotatebox{90}{Precision} &
\rotatebox{90}{Recall} &
\rotatebox{90}{F1-Score} &
\rotatebox{90}{Accuracy} \\
\midrule
\endhead
\bottomrule
\endfoot
 
{\footnotesize Almodovar et al.\ (2024)} & x & x &   &   &   & x & x & x &   &   &   &   &   & -- & -- & {\footnotesize 91.22\%} & -- \\
{\footnotesize Ma et al.\ (2024)}        & x &   & x &   &   & x & x &   &   &   & x &   &   & -- & -- & {\footnotesize 95.7\%} & -- \\
{\footnotesize Liu et al.\ (2024)}       &   &   & x &   &   & x & x &   &   &   &   &   &   & -- & -- & -- & -- \\
{\footnotesize Karlsen et al.\ (2023)}   & x & x &   &   &   &   & x & x &   & x &   &   &   & -- & -- & {\footnotesize 99.8\%} & -- \\
{\footnotesize Wurzenberger et al.\ (2024)} & & & & x &   &   &   &   &   &   &   &   &   & -- & -- & -- & -- \\
{\footnotesize Yang et al.\ (2024)}      &   & x &   & x &   & x & x & x &   &   &   &   &   & -- & -- & {\footnotesize 95.9\%} & -- \\
{\footnotesize Astekin et al.\ (2024a)}  &   &   & x &   &   &   &   &   &   &   & x &   &   & -- & -- & -- & -- \\
{\footnotesize Astekin et al.\ (2024b)}  &   &   & x &   &   &   &   &   &   &   & x &   &   & -- & -- & -- & -- \\
{\footnotesize Partovian et al.\ (2024)} &   &   &   & x &   &   &   &   &   &   &   &   &   & -- & -- & -- & -- \\
{\footnotesize Pham \& Lee (2024)}       &   & x &   &   &   & x & x & x &   &   &   &   &   & -- & -- & {\footnotesize 98.5\%} & -- \\
{\footnotesize Chen et al.\ (2024)}      & x &   & x &   &   &   &   &   &   &   & x &   &   & -- & -- & {\footnotesize 95.3\%} & -- \\
{\footnotesize Krebs et al.\ (2024)}     &   &   &   & x &   &   &   &   & x &   &   &   & x & -- & -- & -- & {\footnotesize 82.84\%} \\
{\footnotesize Mahmoud et al.\ (2024)}   &   &   &   &   &   &   &   &   &   &   &   &   & x & -- & -- & -- & -- \\
{\footnotesize Xu et al.\ (2024a)}       &   &   & x &   &   &   &   &   &   &   & x &   &   & -- & -- & -- & -- \\
{\footnotesize Liao et al.\ (2025)}      &   & x &   &   &   & x & x & x &   &   &   &   &   & -- & -- & {\footnotesize 97.8\%} & -- \\
{\footnotesize Tian et al.\ (2024)}      & x & x &   &   &   & x &   &   &   &   &   &   &   & -- & -- & {\footnotesize 99\%} & -- \\
{\footnotesize Xu et al.\ (2024b)}       &   & x &   &   &   & x & x &   &   &   &   &   &   & -- & -- & {\footnotesize 98.4\%} & -- \\
{\footnotesize Liu et al.\ (2024)}       & x & x &   &   &   &   & x & x &   &   &   &   &   & -- & -- & -- & -- \\
{\footnotesize Ali et al.\ (2025a)}      &   &   &   & x &   & x & x &   &   &   &   &   &   & -- & -- & {\footnotesize 99.78\%} & -- \\
{\footnotesize García Gómez et al.\ (2026)} & & &  & x &   & x & x &   &   &   &   &   &   & -- & -- & {\footnotesize 95\%} & -- \\
{\footnotesize Bhanage \& Pawar (2023)}  &   & x &   & x &   &   &   &   &   &   &   &   &   & -- & -- & -- & {\footnotesize 88.54\%} \\
{\footnotesize Sun et al.\ (2024a)}      &   &   & x &   &   & x & x & x &   &   & x &   &   & -- & -- & -- & -- \\
{\footnotesize Senevirathne et al.\ (2024)} & & & x &  &  &   &   &   &   &   &   &   &   & -- & -- & {\footnotesize 94.5\%} & -- \\
{\footnotesize Aziz \& Munir (2024)}     &   & x &   &   &   & x & x &   &   &   &   &   &   & -- & -- & {\footnotesize 94.5\%} & -- \\
{\footnotesize Corbelle et al.\ (2024)}  &   & x &   &   &   & x &   &   &   &   & x &   &   & -- & -- & -- & -- \\
{\footnotesize Sun et al.\ (2024b)}      &   &   & x &   &   &   &   &   &   &   & x &   &   & -- & -- & -- & -- \\
\midrule
\textbf{Total (26)} & \textbf{6} & \textbf{11} & \textbf{9} & \textbf{6} & \textbf{0} & \textbf{12} & \textbf{12} & \textbf{7} & \textbf{1} & \textbf{1} & \textbf{6} & \textbf{0} & \textbf{2} & -- & -- & -- & -- \\
\end{longtable}
 
%--------------------
% Limpíeza de página
\clearpage
%--------------------

 
% ============================================================
% Tabla 7.2 — Hallazgos parte 2 (papers 27–52)
% ============================================================
\begin{longtable}{p{2.4cm}*{5}{>{\centering\arraybackslash\small}p{0.4cm}}*{5}{>{\centering\arraybackslash\small}p{0.4cm}}*{3}{>{\centering\arraybackslash\small}p{0.4cm}}*{4}{>{\centering\arraybackslash\small}p{0.6cm}}}
 
\caption{Hallazgos por categoría de análisis en artículos seleccionados
         (trabajos 27--52 de 52).}
\label{tab:hallazgos_parte2} \\
 
\toprule
& \multicolumn{5}{c}{\textbf{Modelos / Técnicas}} 
& \multicolumn{5}{c}{\textbf{Datasets}}
& \multicolumn{3}{c}{\textbf{Preproc./KB}}
& \multicolumn{4}{c}{\textbf{Métricas}} \\
\cmidrule(lr){2-6}\cmidrule(lr){7-11}\cmidrule(lr){12-14}\cmidrule(lr){15-18}
\textbf{ID} &
\rotatebox{90}{Fine-tuning} &
\rotatebox{90}{Emb. BERT} &
\rotatebox{90}{LLM/Prompting} &
\rotatebox{90}{ML Clásico} &
\rotatebox{90}{Agente/ReAct} &
\rotatebox{90}{HDFS} &
\rotatebox{90}{BGL} &
\rotatebox{90}{Thunderbird} &
\rotatebox{90}{UWF-Zeek} &
\rotatebox{90}{Logs Web/Red} &
\rotatebox{90}{Log Parsing} &
\rotatebox{90}{RAG / Base KB} &
\rotatebox{90}{MITRE ATT\&CK} &
\rotatebox{90}{Precision} &
\rotatebox{90}{Recall} &
\rotatebox{90}{F1-Score} &
\rotatebox{90}{Accuracy} \\
\midrule
\endfirsthead
\multicolumn{18}{c}{\tablename~\thetable{} ( continuación)} \\
\toprule
& \multicolumn{5}{c}{\textbf{Modelos / Técnicas}} 
& \multicolumn{5}{c}{\textbf{Datasets}}
& \multicolumn{3}{c}{\textbf{Preproc./KB}}
& \multicolumn{4}{c}{\textbf{Métricas}} \\
\cmidrule(lr){2-6}\cmidrule(lr){7-11}\cmidrule(lr){12-14}\cmidrule(lr){15-18}
\textbf{ID} &
\rotatebox{90}{Fine-tuning} &
\rotatebox{90}{Emb. BERT} &
\rotatebox{90}{LLM/Prompting} &
\rotatebox{90}{ML Clásico} &
\rotatebox{90}{Agente/ReAct} &
\rotatebox{90}{HDFS} &
\rotatebox{90}{BGL} &
\rotatebox{90}{Thunderbird} &
\rotatebox{90}{UWF-Zeek} &
\rotatebox{90}{Logs Web/Red} &
\rotatebox{90}{Log Parsing} &
\rotatebox{90}{RAG / Base KB} &
\rotatebox{90}{MITRE ATT\&CK} &
\rotatebox{90}{Precision} &
\rotatebox{90}{Recall} &
\rotatebox{90}{F1-Score} &
\rotatebox{90}{Accuracy} \\
\midrule
\endhead
\bottomrule
\endfoot
 
{\footnotesize Huang et al.\ (2025)}     & x &   & x &   &   & x & x & x &   &   & x &   &   & -- & -- & -- & -- \\
{\footnotesize Boffa et al.\ (2024)}     & x & x &   &   &   &   &   &   &   & x &   &   & x & -- & -- & -- & {\footnotesize 91.2\%} \\
{\footnotesize Zhao et al.\ (2025)}      &   &   & x &   &   &   &   &   &   &   &   &   &   & -- & -- & {\footnotesize 78.8\%} & -- \\
{\footnotesize Palma et al.\ (2025)}     &   &   & x &   &   &   &   &   &   &   &   &   &   & {\footnotesize 89.1\%} & {\footnotesize 90.5\%} & {\footnotesize 92.8\%} & {\footnotesize 97.87\%} \\
{\footnotesize Almorjan et al.\ (2025)}  & x &   &   &   &   &   &   &   &   &   &   &   &   & -- & -- & {\footnotesize 77.43\%} & -- \\
{\footnotesize Cai et al.\ (2025)}       &   & x & x &   &   &   &   &   &   &   &   & x & x & -- & -- & {\footnotesize 98.85\%} & -- \\
{\footnotesize Wang et al.\ (2024)}      &   & x & x &   &   &   &   &   &   &   &   & x &   & -- & -- & -- & -- \\
{\footnotesize Nasution et al.\ (2025)}  &   &   & x &   &   &   &   &   &   &   &   &   &   & -- & -- & -- & -- \\
{\footnotesize Belcastro et al.\ (2026)} &   & x & x &   &   &   &   &   &   &   &   & x &   & -- & -- & {\footnotesize 84.1\%} & -- \\
{\footnotesize Blefari et al.\ (2026)}   & x & x & x &   & x &   &   &   &   & x &   & x & x & {\footnotesize 97\%} & {\footnotesize 97\%} & {\footnotesize 97\%} & -- \\
{\footnotesize Lee et al.\ (2025)}       &   &   & x &   & x &   &   &   &   &   &   &   & x & {\footnotesize 97\%} & {\footnotesize 97\%} & {\footnotesize 97\%} & {\footnotesize 99\%} \\
{\footnotesize Kshetri (2025)}           &   &   &   &   & x &   &   &   &   &   &   & x &   & -- & -- & -- & -- \\
{\footnotesize Ali et al.\ (2025b)}      &   &   &   & x &   &   &   &   &   &   &   &   &   & -- & -- & -- & {\footnotesize 99.6\%} \\
{\footnotesize Ismail et al.\ (2025)}    &   &   & x &   & x &   &   &   &   &   &   & x & x & -- & -- & -- & -- \\
{\footnotesize Chamkar et al.\ (2025)}   &   &   &   &   &   &   &   &   &   &   &   &   & x & -- & -- & -- & -- \\
{\footnotesize Pitkar (2025)}            &   &   &   &   &   &   &   &   &   &   &   &   &   & -- & -- & -- & -- \\
{\footnotesize Ye et al.\ (2025)}        &   &   &   & x &   &   &   &   &   &   &   &   &   & -- & -- & {\footnotesize 99.9\%} & -- \\
{\footnotesize Hmimou et al.\ (2025)}    &   &   & x &   & x &   &   &   &   &   &   & x &   & -- & -- & -- & {\footnotesize 93.6\%} \\
{\footnotesize Binbeshr et al.\ (2025)}  &   &   &   &   &   &   &   &   &   &   &   &   &   & -- & -- & -- & -- \\
{\footnotesize Falowo et al.\ (2024)}    &   &   &   &   & x &   &   &   &   &   &   &   &   & -- & -- & -- & -- \\
{\footnotesize Kapera \& Niemiec (2024)} &   &   &   & x &   &   &   &   &   &   &   &   &   & -- & -- & -- & -- \\
{\footnotesize Rigaki et al.\ (2023)}    &   &   & x &   & x &   &   &   &   &   &   &   &   & -- & -- & -- & {\footnotesize 83.3\%} \\
{\footnotesize Ahmed et al.\ (2024)}     &   &   & x &   &   &   &   &   &   &   &   &   &   & -- & -- & -- & -- \\
{\footnotesize Shah \& Madisetti (2025)} &   &   & x &   & x &   &   &   &   &   &   & x &   & -- & -- & {\footnotesize 67.82\%} & -- \\
{\footnotesize Yuan et al.\ (2024)}      &   &   & x &   & x &   &   &   &   &   &   &   &   & -- & -- & -- & {\footnotesize 82.7\%} \\
{\footnotesize He et al.\ (2023)}        &   & x &   &   &   & x & x &   &   &   &   &   &   & -- & -- & {\footnotesize 97.1\%} & -- \\
\midrule
\textbf{Total (26)} & \textbf{4} & \textbf{6} & \textbf{15} & \textbf{4} & \textbf{9} & \textbf{3} & \textbf{3} & \textbf{1} & \textbf{0} & \textbf{2} & \textbf{3} & \textbf{8} & \textbf{6} & -- & -- & -- & -- \\
\midrule
\textbf{Total general (52)} & \textbf{10} & \textbf{17} & \textbf{24} & \textbf{10} & \textbf{9} & \textbf{15} & \textbf{15} & \textbf{8} & \textbf{1} & \textbf{3} & \textbf{9} & \textbf{8} & \textbf{8} & -- & -- & -- & -- \\
\end{longtable}


Los resultados de la Tabla~\ref{tab:hallazgos_parte1} y la
Tabla~\ref{tab:hallazgos_parte2} revelan que \textit{LLM/Prompting}
es la técnica más frecuente en el corpus (24 de 52 trabajos),
seguida por \textit{Embeddings BERT} (17 trabajos), lo que confirma
la tendencia hacia representaciones semánticas profundas en el área.
En cuanto a datasets, HDFS y BGL son los benchmarks de referencia
predominantes (15 trabajos cada uno), siendo Thunderbird el tercero
más usado (8 trabajos). Respecto al preprocesamiento e integración
de conocimiento, solo 8 trabajos incorporan bases RAG y 8 integran
el framework MITRE ATT\&CK, confirmando que la combinación de
razonamiento semántico con inteligencia de amenazas estructurada
permanece como espacio abierto. Las métricas más reportadas son
F1-Score y Accuracy, con rangos que van del 67.82\,\% al 99.90\,\%
en F1 dependiendo de la complejidad de la tarea y el dominio de logs.
 


%%-------------------------------------------------------------
\section{RQ1: Técnicas para el análisis de logs en sistemas
  complejos}
%%-------------------------------------------------------------
Esta sección sintetiza las técnicas identificadas en el corpus para
el análisis automatizado de logs en sistemas complejos. Los trabajos
revisados abordan dos dimensiones complementarias: el \textit{parsing}
de logs y la detección de anomalías. Ambas
dimensiones han evolucionado desde enfoques estadísticos clásicos
hacia modelos semánticos basados en LLMs, con una tendencia creciente
hacia la integración de razonamiento contextual y frameworks de
evaluación reproducibles.

\subsection{Del parsing estadístico a los embeddings semánticos}

El análisis automatizado de logs de seguridad ha seguido una
trayectoria evolutiva clara en la literatura reciente. Los
enfoques tradicionales se apoyan en parsers estadísticos como
Drain, que extraen plantillas de log mediante agrupación por
similitud léxica. Si bien Drain ha sido el estándar de facto
durante años, el estudio comparativo de Astekin et
al.~\cite{astekin2024comparative} cuantifica su limitación:
CodeLlama ejecutado localmente vía Ollama supera a GPT-3.5 en
Parsing Accuracy (PA) con prompting few-shot, validando que
los modelos open-source son alternativas viables sin costo de
API. El estudio sobre no-determinismo del mismo equipo~\cite{astekin2024nondeterminism}
aporta un hallazgo crítico de ingeniería: temperatura cero
reduce pero no elimina la variabilidad en salidas de LLMs para
parsing, identificando que los modelos de código (CodeLlama)
son sensiblemente más estables, lo que informa el diseño del
agente propuesto en este trabajo.

LLMParser~\cite{ma2024llmparser} reformula el parsing como
una tarea de traducción texto-a-texto, demostrando que
Flan-T5-base (240M parámetros) alcanza un PA de 0.957 con
\textit{few-shot tuning}, superando a Drain en un factor de
$2.85\times$ en precisión. Hooglle~\cite{chen2024hooglle}
da un paso adicional al proponer un parser online de dos fases
que combina un LLM fine-tuneado (p-tuning v2) con un árbol de
prefijos a nivel de carácter (CI-tree), logrando PA$=0.953$
y GA$=0.987$ a velocidad comparable a los parsers estadísticos
, donde el 99.98\,\% de los logs se asignan por coincidencia exacta
sin invocar al LLM. DivLog~\cite{xu2024divlog} lleva esta
dirección al extremo con un parser completamente \textit{training-free}
basado en In-Context Learning (ICL) sobre GPT-3, estableciendo
que es posible obtener calidad de parsing competitiva sin
ningún dato etiquetado.

Una dimensión frecuentemente ignorada en los parsers es el
contenido semántico de las partes variables de los logs. El
trabajo de Wurzenberger et al.~\cite{wurzenberger2024vtd}
aborda este problema mediante el Variable Type Detector (VTD),
un enfoque no supervisado que clasifica estadísticamente cada
token variable mediante tests Kolmogorov-Smirnov, Durbin-Watson
y $\chi^2$, superando a PCA y análisis de series temporales en
curvas ROC sobre los 8 \textit{testbeds} del dataset AIT-LDS v2. VISTA~\cite{sun2024variables}
complementa esta perspectiva con una estrategia centrada en variables
que reorganiza el pipeline de parsing para que las partes dinámicas
fluyan como dato principal, logrando un \textit{speedup} de $2.3\times$ sobre
LILAC en tiempo de muestreo. LogRules~\cite{huang2025logrules}
aporta un enfoque diferente: genera reglas explícitas con GPT-4o-mini
que guían a LLMs de 8B parámetros (LLaMA~3) en el análisis de logs,
logrando rendimiento comparable a modelos propietarios más grandes.



La Figura~\ref{fig:comparativa_parsers} compara la Parsing Accuracy (PA)
y Grouping Accuracy (GA) de los principales parsers de logs identificados
en el corpus. LLMParser y Hooglle superan el baseline de Drain (PA$=0.856$)
con enfoques basados en \textit{fine-tuning}, mientras que DivLog, VISTA y LogRules
no reportan métricas numéricas de PA/GA al emplear estrategias cualitativas
o evaluaciones propias de cada dataset.

\begin{figure}[H]
\centering
\includegraphics[width=1\textwidth]{Resource/rq1-1.png}
\caption{Comparativa de Parsing Accuracy (PA) y Grouping Accuracy (GA)
         de los principales técnicas de \textit{log parsing} identificadas
         en el corpus SLR. La línea discontinua indica el baseline de Drain
         (PA$=0.856$). Fuente: Elaboración propia a partir de datos del corpus.}
\label{fig:comparativa_parsers}
\end{figure}

\subsection{Detección de anomalías: de ML clásico a modelos semánticos}

La detección de anomalías en logs ha transitado por tres generaciones
de técnicas, documentadas comprensivamente por el estudio empírico de
Ali et al.~\cite{ali2025mlstudy}, que evalúa 9 métodos (SVM, Random Forest,
LSTM, LogRobust, NeuralLog, DeepLog, entre otros) sobre 7 datasets de
referencia. El hallazgo central es que los métodos supervisados
tradicionales como Random Forest alcanzan F1$=99.78$ en HDFS, mientras
que DeepLog colapsa a F1$=0.00$ ante el \textit{challenge} de variables en
mensajes, evidenciando que la robustez semántica es una dimensión
independiente de la precisión nominal.

SemPCA~\cite{yang2024sempca} demuestra que añadir representación
semántica (TF-IDF + GloVe) a un algoritmo de Análisis de Componentes Principales (PCA) clásico permite igualar
el rendimiento de los mejores métodos de deep learning supervisados,
con un F1 promedio de 0.959, siendo al menos $5.8\times$ más rápido
en entrenamiento. MultiLog~\cite{pham2024multilog} lleva esta idea más
lejos mediante multi-task learning semi-supervisado con tres objetivos
simultáneos (predicción del próximo template, predicción de templates
enmascarados y minimización de distancia al centroide normal), obteniendo
un modelo de solo 204.290 parámetros entrenables con F1 promedio de 0.985
sobre HDFS, BGL y Thunderbird.

LogBASA~\cite{liao2025logbasa} propone un grafo de conocimiento de comportamiento
del sistema (SLKG) combinado con awareness semántico global, obteniendo
mejoras de F1 sobre 8 baselines en HDFS, BGL y Thunderbird. CLDTLog~\cite{tian2024cldtlog}
introduce aprendizaje contrastivo con tareas de doble objetivo sobre BERT,
logrando F1$=0.99$ en HDFS sin necesidad de log parsing previo. LogCEM~\cite{xu2024semantic}
combina CNN multiscala, Efficient Channel Attention y Mogrifier GRU con
\textit{embeddings} BERT + TF-IDF, alcanzando F1$=98.40\,\%$ en HDFS y
$98.34\,\%$ en BGL.

El Systematic Mapping Study de Partovian et al.~\cite{partovian2024sms}
analiza 74 estudios primarios sobre análisis de logs en entornos
Industry~4.0, concluyendo que la «incapacidad de identificar la causa raíz del problema» es el desafío de capacidad más frecuente en la literatura identificando el razonamiento causal como brecha abierta que este seminario busca cerrar.

LogBD~\cite{liu2024logbd} aborda el problema cross-sistema mediante
adaptación de dominio adversarial sobre BERT, logrando que un modelo
entrenado en Thunderbird se transfiera a BGL con AUC superior a los
métodos sin adaptación. SemiSMAC~\cite{sun2024semismac} propone un
framework semi-supervisado donde ChatGPT realiza simultáneamente el
parseo y el agrupamiento semántico, con hiper-parámetros optimizados
automáticamente. Bhanage \& Pawar~\cite{bhanage2023bert} establecen
que \textit{embeddings} BERT combinados con clasificadores clásicos (SVM, GBDT)
alcanzan accuracy sobre $88\,\%$ en logs Apache sin requerir
arquitecturas de deep learning complejas.


La Figura~\ref{fig:f1_anomaly} compara el F1-Score reportado por
los principales métodos de detección de anomalías en logs del
corpus. La comparación evidencia la brecha entre métodos clásicos
supervisados (con F1 próximo a 1.0 en condiciones
controladas) y métodos semánticos, cuya ventaja radica en la
robustez ante variabilidad léxica y escenarios cross-sistema donde
los enfoques estadísticos colapsan.

\begin{figure}[H]
\centering
\includegraphics[width=1\textwidth]{Resource/rq1-2.png}
\caption{F1-Score de los principales m\'{e}todos de detecci\'{o}n de anomal\'{i}as
         en logs reportados en el corpus del SLR. El valor de DeepLog
         corresponde al \textit{Challenge~3} (variables en mensajes)
         reportado por Garc\'{i}a G\'{o}mez et al.~\cite{garcia2026collaborative}.
         Fuente: Elaboraci\'{o}n propia a partir de datos del corpus.}
\label{fig:f1_anomaly}
\end{figure}

\subsection{Infraestructura, datasets y frameworks de evaluación}

El trabajo de Mahmoud et al.~\cite{mahmoud2024aau} contribuye al área
mediante la construcción del dataset público AAU\_MalData, generado
sobre un entorno de sandbox con Sysmon y ELK Stack, cubriendo
2.800 muestras de malware (89 familias) correlacionadas con MITRE
ATT\&CK. Krebs et al.~\cite{krebs2024svm} aplican clasificación
multi-clase SVM sobre el dataset UWF-ZeekDataFall22 (700.340 logs de
red Zeek con tácticas MITRE), obteniendo \textit{accuracy} del 82.84\% con
OvO pero con desbalance severo del 61\% en una sola táctica,
ilustrando cuantitativamente por qué los métodos clásicos sin
comprensión semántica son insuficientes.

En García Gómez et al.~\cite{garcia2026collaborative} aportan una taxonomía
de métodos LAD en tres categorías (Sequential, Embedding, Graph-wise)
junto con el framework open-source \texttt{log-gym} para evaluación
reproducible. Su hallazgo sobre el fallo de DeepLog en Challenge~3
(variables en mensajes de log) es directamente citable como motivación
del enfoque semántico del seminario. Senevirathne et al.~\cite{senevirathne2024vm}
extienden el análisis de anomalías al dominio de tolerancia a fallos
en VMs, combinando Matrix Profile con inferencia LLM para desencadenar
migración proactiva.

AnyLog~\cite{aziz2024deeplog} propone un pipeline híbrido no supervisado
(BERT + Self-Organizing Maps + Autoencoder) que alcanza F1$=94.5\%$
en \textit{training} y $91\%$ en \textit{testing} sobre conjuntos combinados HDFS + BGL
sin depender de plantillas predefinidas. HC-BERT~\cite{corbelle2024semantic}
aporta una perspectiva de compactación semántica de logs usando
clasificación jerárquica con WordNet, logrando reducir 2.000 líneas de
HDFS a solo 16 líneas compactadas con 99.20\% de reducción preservando
la semántica del evento.

LogFiT~\cite{almodovar2024logfit} cierra este panorama con un enfoque
que elimina el paso de parsing: fine-tunea RoBERTa y Longformer sobre
datos normales únicamente mediante Masked Sentence Prediction. Alcanza
F1$=91.22$ en HDFS y, más importante, F1$=89.38$ en el escenario con
variabilidad léxica (BGL con sinónimos WordNet), donde DeepLog colapsa
a 53.38, demostrando resistencia al \textit{log drift}.

%%-------------------------------------------------------------
\section{RQ2: Uso de LLMs para análisis de logs en ciberseguridad}
%%-------------------------------------------------------------

Esta sección analiza el uso de modelos de lenguaje de gran escala
para el análisis semántico de logs en contextos de ciberseguridad.
Los trabajos revisados abordan tres dimensiones: el diseño de
estrategias de \textit{prompting }para tareas \textit{zero-shot} y \textit{few-shot}, la
adaptación mediante \textit{fine-tuning} a dominios específicos de seguridad,
y la extracción de conocimiento estructurado para alimentar bases de
inteligencia de amenazas. En conjunto, el corpus evidencia una
transición desde la clasificación supervisada hacia arquitecturas
agénticas que integran LLMs con bases de conocimiento como MITRE
ATT\&CK para razonamiento explicable.

\subsection{Prompt engineering y razonamiento zero-shot}

La investigación de Liu et al.~\cite{liu2024logprompt}, que propone
LogPrompt, es el antecedente más directo en interpretabilidad y
razonamiento explícito aplicado a logs. LogPrompt aplica tres
estrategias de prompting: self-prompt, Chain-of-Thought (CoT)
implícito/explícito e in-context learning sobre ChatGPT en escenario
completamente zero-shot. En detección de anomalías, CoT explícito
mejora en 62.4\,\% sobre el prompt simple en BGL y en 380.7\,\% sobre
Vicuna~13B. La evaluación de interpretabilidad con 6 expertos asigna
un HIP (High Interpretability Percentage) del 83\,\% en parsing y
84.5\,\% en \textit{anomaly detection}, estableciendo criterios formales de
calidad narrativa que el presente trabajo adopta para su evaluación.

TAD-GP~\cite{zhao2025tadgp} demuestra que modelos de 7B parámetros
pueden superar a GPT-3.5-turbo en tareas de ciberseguridad cuando se
aplica prompting estructurado. Usando Gemma-7b-it con CoT en cinco
pasos y diálogo multi-turno virtual, logra F1 promedio de 0.788 sobre
KDD~Cup~1999, CICIDS2017 y UNSW-NB15, frente a 0.715 de GPT-3.5-turbo.
Este resultado evidencia que la calidad del razonamiento depende
críticamente del diseño del prompt y no solo de la escala del modelo.

Palma et al.~\cite{palma2025llm} realizan la evaluación más alineada
con la arquitectura del presente seminario: comparan seis LLMs locales
vía Ollama (incluyendo qwen2.5:7b, llama3.2:3b y gemma3:4b) en
clasificación de logs de seguridad empresariales reales. El mejor
modelo, qwen2.5:7b, alcanza accuracy del 97.87\,\% y F1$=0.928$, con
tiempo promedio de 3.3 segundos por caso. El paper demuestra
empíricamente que el \textit{role-anchoring} en prompts combinado con
formato estructurado obligatorio mejora la consistencia de las
respuestas.

Nasution et al.~\cite{nasution2025phishing} realizan el primer benchmark
de 21 LLMs open-source para detección de \textit{phishing} bajo cuatro estrategias
de \textit{prompting}, encontrando que los modelos más grandes con few-shot
superan consistentemente a los zero-shot, con variaciones de F1 de
hasta 20 puntos porcentuales según la estrategia. El estudio empírico de
SemiSMAC~\cite{sun2024semismac} confirma que ChatGPT con un solo paso
de parseo y agrupamiento semántico supera a Drain en Grouping Accuracy
sobre múltiples datasets, a costo de mayor latencia.

La Figura~\ref{fig:f1_llm_ciberseguridad} compara el F1-Score
reportado por los principales sistemas LLM para análisis de logs
en ciberseguridad identificados en el corpus. Se observa que los
sistemas con arquitecturas agénticas y RAG (CyberRAG, LogRESP-Agent)
alcanzan los mejores resultados, mientras que los enfoques de
prompting puro sin integración de bases de conocimiento presentan
menor rendimiento.

\begin{figure}[H]
\centering
\includegraphics[width=1\textwidth]{Resource/rq2-1.png}
\caption{F1-Score de los principales sistemas basados en LLMs para
         análisis de logs en ciberseguridad reportados en el corpus SLR.
         Fuente: Elaboración propia a partir de datos del corpus.}
\label{fig:f1_llm_ciberseguridad}
\end{figure}

\subsection{Fine-tuning y adaptación al dominio}

LLM4Sec~\cite{karlsen2023llm4sec} provee el benchmark de referencia
para fine-tuning supervisado en análisis de logs de seguridad: evalúa
5 arquitecturas sobre 6 datasets combinando logs web (Apache, CSIC, PKDD)
y de sistema (Thunderbird, Spirit, BGL), reportando que DistilRoBERTa
fine-tuneada alcanza F1 promedio de 0.998. El módulo de interpretabilidad
TranSHAP identifica correctamente las características de ataque a nivel
de token, estableciendo precedente de XAI para logs de seguridad.

LogPrécis~\cite{boffa2024logprecis} es el antecedente más directo al
seminario: recibe sesiones de shell Unix maliciosas de honeypot y asigna
tácticas MITRE ATT\&CK mediante NER fine-tuneada sobre CodeBERT con solo
360 sesiones etiquetadas. La configuración ganadora alcanza
\textit{accuracy}$=0.912$ y \textit{Fidelity}$=0.669$, superando a GPT-3 Davinci con
tiempo de inferencia $23\times$ menor. La reducción de 400.000 sesiones
a 3.000 fingerprints distintos habilita análisis forense a escala. Las
diferencias con el seminario son: LogPrécis opera sobre un único tipo de
log, clasifica tácticas token a token sin reconstruir la narrativa causal,
y requiere etiquetado supervisado.

CyberRAG~\cite{blefari2026cyberrag} implementa un framework modular de
Agentic RAG usando LLaMA~3.1:8B vía Ollama para clasificación de ataques
web. El pipeline incluye clasificadores BERT especializados por familia
de ataque, recuperación iterativa MMR desde bases de conocimiento
(PortSwigger, OWASP, CVE/NVD, MITRE ATT\&CK), y un módulo de generación
de reportes en lenguaje natural. Las métricas F1$=0.97$ (binario) y
F1$=0.99$ (multi-clase) son estadísticamente significativas (Welch's
t-test, $p<0.001$).

LogRESP-Agent~\cite{lee2025logrespagent} extiende la idea agentica con
un ciclo recursivo Thought-Action-Observation (TAO): el agente detecta
anomalías, mapea TTPs a MITRE ATT\&CK, genera explicaciones auditables
e itera con nuevas observaciones. Usando Gemini-2.0-Flash con 7 herramientas
modulares, alcanza F1$=0.97$ en detección binaria y $F1=0.99$ en
clasificación multi-clase sobre el dataset Monster-THC.


La Figura~\ref{fig:comparativa_agentes} compara los sistemas
agénticos más relevantes del corpus según sus métricas de
F1 binario, F1 multi-clase y Accuracy. CyberRAG y LogRESP-Agent
lideran en ambas métricas de F1, validando la efectividad del
patrón Agentic RAG con integración de MITRE ATT\&CK para
clasificación de incidentes de seguridad.

\begin{figure}[H]
\centering
\includegraphics[width=1\textwidth]{Resource/rq2-2.png}
\caption{Comparativa de sistemas agénticos LLM para SOC según
         F1 binario, F1 multi-clase y Accuracy reportados en el
         corpus SLR. Fuente: Elaboración propia a partir de datos
         del corpus.}
\label{fig:comparativa_agentes}
\end{figure}

\subsection{Extracción de conocimiento y threat intelligence}

Almorjan et al.~\cite{almorjan2025ioc} demuestran que los LLMs pueden
generar datos sintéticos de alta calidad para Indicadores de Compromiso
(IoCs) mediante fine-tuning de GPT-3.5-turbo. Los modelos DL entrenados
sobre datos sintéticos alcanzan F1$=77.43\,\%$ en clasificación de IoCs,
con clases como SHA-1, SHA-256 y dominios con precisión $\geq 97\,\%$.

ESP~\cite{cai2025esp} aporta la evidencia más sólida sobre extracción
de conocimiento de reportes de amenazas (CTRs) para alimentar bases RAG:
combina BERT + BiLSTM + Dual-Sequence Attention para NER y Qwen2.5-72B
mediante ICL para extracción de relaciones, alcanzando
F1$=98.85\,\%$ sobre Enterprise\_RE. La generación automática de reglas
Sigma con referencias a técnicas ATT\&CK (\textit{e.g.}, \texttt{attack.t1057})
es análoga al mapeo que el agente del seminario realiza entre eventos de
log y técnicas MITRE. El paper valida además que STIX~2.1 es el esquema ontológico apropiado para representar entidades
de ciberseguridad.

KLAGE~\cite{belcastro2026klage} construye Knowledge Graphs a partir de
flujos de tráfico de red, los clasifica con Graph-BERT (F1$=84.1\,\%$
sobre RT-IoT2022, superando al mejor método GNN previo en $+5$ puntos),
aplica LIME para explicabilidad post-hoc, y usa GPT-4o para generar
reportes narrativos. El ablation study demuestra que el contexto
estructurado es indispensable: GPT-adv (con grafo + LIME) obtiene score
de explicabilidad de 0.80, frente a 0.22 de GPT-base (sin contexto).
Esta evidencia valida directamente la arquitectura del modelo planteado en este trabajo, donde
el agente recibe contexto estructurado antes de generar la narrativa.

MAD-CTI~\cite{shah2025madcti} aporta un framework multi-agente para
extracción y clasificación de CTI desde la dark web, alcanzando F1
ponderado de 0.678 en clasificación de contenido relevante. Yuan et
al.~\cite{yuan2024osint} demuestran que los LLMs con herramientas externas
(patrón agente-con-herramientas) mejoran dramáticamente en adquisición
de OSINT: Qwen-7B-Chat con el framework de agente (82.7\,\% accuracy en
CVEs) supera en $6.2\times$ al mismo modelo sin herramientas (13.3\,\%).

%%-------------------------------------------------------------
\section{RQ3: Desafíos de los SOC en análisis de logs}
%%-------------------------------------------------------------

Esta sección examina los desafíos operacionales que enfrentan los
Centros de Operaciones de Seguridad (SOC) en el análisis de logs y
la respuesta a incidentes, y cómo los trabajos del corpus proponen
soluciones basadas en automatización e inteligencia artificial. Los
estudios revisados abordan tres dimensiones: la sobrecarga de alertas
y la reducción de falsos positivos, el diseño de arquitecturas SOAR
con agentes LLM, y el análisis de brechas en la cobertura de detección
frente a marcos como MITRE ATT\&CK.

\subsection{Sobrecarga de alertas y automatización de la respuesta}

Los Centros de Operaciones de Seguridad (SOC) enfrentan una crisis
estructural documentada: la escasez global de 4 millones de profesionales
de ciberseguridad coexiste con \textit{breakout times} de ataques de apenas 2
minutos~\cite{kshetri2025agentic}. Kshetri~\cite{kshetri2025agentic}
sistematiza el impacto de la IA agentiva mediante casos industriales:
Charlotte AI (CrowdStrike) reporta 98\,\% de precisión en \textit{triage}
eliminando más de 40 horas semanales de trabajo manual; GreyMatter
(ReliaQuest) procesa alertas $20\times$ más rápido y automatiza el
98\,\% con contención en menos de 5 minutos; Microsoft Security Copilot
reduce el MTTR en 30\,\% y hace a analistas novatos 26\,\% más rápidos.

La revisión sistemática de Binbeshr et al.~\cite{binbeshr2025socsrl}
analiza 38 estudios (2006--2024) sobre IA en SOC cognitivos, identificando
que el crecimiento de publicaciones alcanzó su pico en 2021 y que las
principales oportunidades abiertas son: razonamiento explicable, integración
multi-fuente, y reducción de dependencia de expertos humanos.

Falowo et al.~\cite{falowo2024idr} formalizan matemáticamente la
diferencia entre SOC basados en AI tradicional y SOC hipotéticos
basados en AGI, estableciendo métricas de TP/FP formales para incidentes.
Kapera \& Niemiec~\cite{kapera2024siem} aportan evidencia cuantitativa
de la problemática de falsos positivos en SIEM: su sistema DRBA reduce
el promedio diario de falsos positivos de 8.74 a 6.42 (reducción del
26.6\,\%) sobre 31 días de producción en Splunk, con rango de mejora
del 7.8--45.4\,\% según período.

El trabajo de Ali et al.~\cite{ali2025apt} aborda la detección de APTs
mediante Splunk como plataforma SIEM, con seis módulos de detección por
categoría y XGBoost como clasificador final. El sistema alcanza accuracy
del 99.6\,\% y FPR$=0.0\,\%$, superando al mejor trabajo previo
(84.8\,\% con FPR$=4.5\,\%$). Sin embargo, entrega únicamente una
clasificación binaria sin razonamiento causal ni generación de narrativas
explicables.


La Figura~\ref{fig:impacto_ia_soc} resume las métricas de impacto
reportadas por los trabajos del corpus que evalúan sistemas de
automatización IA en operaciones SOC. Los valores evidencian que
la integración de IA reduce significativamente la carga operacional,
con mejoras que van desde el 26.6\,\% en reducción de falsos positivos
hasta el 99.6\,\% de accuracy en detección de APTs.

\begin{figure}[H]
\centering
\includegraphics[width=1\textwidth]{Resource/rq3-1.png}
\caption{Métricas de impacto de la automatización con IA en operaciones
         SOC reportadas en el corpus SLR. Fuente: Elaboración propia
         a partir de datos del corpus.}
\label{fig:impacto_ia_soc}
\end{figure}

\subsection{Arquitecturas SOAR y agentes LLM para SOC}

El framework IVAM~\cite{ismail2025ivam} propone un agente LLM autónomo
integrado con Wazuh (SIEM open-source), MITRE ATT\&CK como base de
conocimiento para mapeo de TTPs, RAG sobre historial de incidentes
previos, y generación dinámica de pasos técnicos de respuesta. El
agente condensa el playbook estándar de 38 pasos de Palo Alto XSOAR
para \textit{brute-force} SSH en 10 acciones equivalentes, y el QRA demuestra
reducción medible del riesgo relativo.

El sistema de Hmimou et al.~\cite{hmimou2025multiagent} implementa cuatro
agentes especializados coordinados por un Cross-Context Recommendation
System, reportando accuracy del 93.6\,\%, reducción de falsos positivos
del 41.3\,\% frente a un SIEM baseline, y triage 38.5\,\% más rápido.
El Log Analyzer Agent combina ELK Stack, Suricata y
LLaMA~3.3-70B para interpretación semántica.

CyberRAG~\cite{blefari2026cyberrag} y LogRESP-Agent~\cite{lee2025logrespagent}
representan el estado del arte más reciente en \textit{pipelines} agenticos para
clasificación de incidentes, como se detalló en la sección anterior.


La Figura~\ref{fig:comparativa_soc_frameworks} compara los principales
frameworks de automatización SOC del corpus según accuracy, reducción
de falsos positivos y eficiencia en ejecución de playbooks. IVAM
destaca por la reducción de pasos de respuesta, mientras que Ali et
al.\ y Hmimou et al.\ lideran en accuracy y reducción de falsos
positivos respectivamente.

\begin{figure}[H]
\centering
\includegraphics[width=1\textwidth]{Resource/rq3-2.png}
\caption{Comparativa de frameworks de automatización para SOC según
         accuracy, reducción de falsos positivos y eficiencia en
         playbooks reportadas en el corpus SLR. Fuente: Elaboración
         propia a partir de datos del corpus.}
\label{fig:comparativa_soc_frameworks}
\end{figure}

\subsection{Gap analysis y cobertura de detección SOC}

Chamkar et al.~\cite{chamkar2025soc} aportan una perspectiva metodológica sobre la forma en que los SOC evalúan su cobertura de detección frente a MITRE ATT\&CK. En particular, el mapeo manual entre Event ID y técnicas MITRE desarrollado en dicho trabajo corresponde al mismo tipo de conocimiento que el RAG sobre MITRE ATT\&CK propuesto en este seminario busca proporcionar de manera automatizada, lo que respalda la pertinencia de esta decisión arquitectónica.

Pitkar~\cite{pitkar2025cloud} proyecta que el mercado de AI en
ciberseguridad crecerá de USD~24.3B en 2023 a USD~134B en 2030 y que
el 99\,\% de las brechas cloud serán atribuibles a errores humanos o
misconfiguraciones. García Gómez et al.~\cite{garcia2026collaborative}
evidencian que los métodos LAD colaborativos distribuidos enfrentan
desafíos de heterogeneidad de datos y latencia que afectan su aplicabilidad
en SOC pequeños.

Ye et al.~\cite{ye2025insider} abordan la detección de amenazas internas
desde el ángulo de la privacidad, alcanzando F1$=0.999$ con Federated
Learning sobre el dataset CERT~v5.2. Mahmoud et al.~\cite{mahmoud2024aau}
validan que el stack Sysmon + ELK es la contraparte open-source de la
infraestructura SIEM que el prototipo propuesto abstrae, proporcionando
un dataset alternativo con 2.800 muestras de malware etiquetadas.

%%-------------------------------------------------------------
\section{RQ4: Capacidades de razonamiento de LLMs en ciberseguridad}
%%-------------------------------------------------------------

Esta sección analiza las capacidades de razonamiento de los modelos
de lenguaje de gran escala aplicadas a tareas de ciberseguridad,
con énfasis en los patrones arquitectónicos que habilitan razonamiento
iterativo, interpretabilidad y correlación multi-fuente. Los trabajos
revisados abordan cuatro dimensiones: el patrón ReAct como referente
para agentes autónomos, el impacto de las estrategias de prompting
sobre el rendimiento, la explicabilidad mediante técnicas XAI, y el
razonamiento causal sobre múltiples fuentes de evidencia.

\subsection{El patrón ReAct como referente arquitectónico}

El trabajo de Rigaki et al.~\cite{rigaki2023cage} constituye un antecedente relevante para el uso de agentes basados en LLM en tareas de ciberseguridad, al demostrar que la incorporación de razonamiento explícito y retroalimentación del entorno mejora el desempeño frente a enfoques de una sola instrucción. En NetSecGame, los agentes LLM son evaluados en escenarios de exfiltración APT mediante variantes Single-Prompt y ReAct. Los resultados muestran que ReAct + GPT-4 alcanza un \textit{win rate} del 100,\,\% en el entorno sin defensor y del 83.3\,\% con defensor estocástico, desempeño comparable al de un agente Q-learning. Más allá del patrón ReAct en sí mismo, la contribución principal para este seminario radica en evidenciar que los agentes que observan el estado del entorno, razonan sobre acciones previas y ajustan sus decisiones evitan repetir acciones fallidas y logran un comportamiento más robusto.

LogRESP-Agent~\cite{lee2025logrespagent} es la implementación más reciente
del patrón agentico aplicado a logs: su ciclo TAO recursivo, con memoria
de observaciones previas y suite de 7 herramientas modulares (TTPMapper,
ThreatLookup, DescriptionGenerator, entre otros), es directamente comparable
al diseño del seminario. Yuan et al.~\cite{yuan2024osint} validan la
arquitectura herramienta-con-agente: Qwen-7B-Chat con acceso a herramientas
OSINT mejora en $6.2\times$ en accuracy sobre CVEs frente a la versión
sin herramientas, cuantificando el valor de la extensión agentica.


La Figura~\ref{fig:capacidades_razonamiento} compara las métricas
de capacidad de razonamiento agéntico y explicabilidad reportadas
en el corpus. La brecha entre KLAGE con contexto estructurado
(80.0) y sin contexto (22.0) cuantifica el valor del razonamiento
fundamentado en evidencia, validando directamente la arquitectura
RAG + MITRE ATT\&CK propuesta en este seminario.

\begin{figure}[H]
\centering
\includegraphics[width=1\textwidth]{Resource/rq4-2.png}
\caption{Capacidades de razonamiento agéntico y explicabilidad (XAI)
         reportadas en el corpus SLR. Fuente: Elaboración propia
         a partir de datos del corpus.}
\label{fig:capacidades_razonamiento}
\end{figure}

\subsection{Estrategias de prompting y razonamiento sobre frameworks de seguridad}

Ahmed et al.~\cite{ahmed2024cis} demuestran que los LLMs pueden razonar
efectivamente sobre controles de seguridad: combinando \textit{few-shot prompting},
CoT y \textit{generated knowledge prompting} en cadena, automatizan la validación
de los 153 controles del framework CIS~CSC~v8. GPT-4 extrae correctamente
Event IDs de Windows y Linux para cada control, generando reglas SIEM
válidas. La evaluación mediante LLM-as-a-judge (correlación
Pearson$=0.60$ con evaluación humana) es directamente aplicable al
componente de evaluación de narrativa del seminario (narrative quality
$\geq 3.0/4.0$).

Nasution et al.~\cite{nasution2025phishing} sistematizan el efecto de la
estrategia de \textit{prompting}: la diferencia entre zero-shot y few-shot CoT
puede alcanzar 20 puntos de F1 en el mismo modelo, indicando que la
estrategia de prompting es al menos tan importante como la elección del
modelo. LogRules~\cite{huang2025logrules} extiende esta idea
al generar reglas explícitas que estructuran el razonamiento del LLM,
demostrando que modelos de 8B parámetros con guía de reglas superan a
modelos propietarios más grandes sin guía.

Kshetri~\cite{kshetri2025agentic} caracteriza conceptualmente los
componentes de agentes autónomos (memoria, planificación, perfil, acción)
y analiza el Agentic RAG, donde el agente refina consultas iterativamente.


La Figura~\ref{fig:impacto_prompting} ilustra el impacto de las
distintas estrategias de prompting sobre el rendimiento de los LLMs
en tareas de ciberseguridad, según los trabajos del corpus. La
progresión desde zero-shot hasta role-anchoring con formato
estructurado evidencia que la estrategia de prompting es al menos
tan determinante como la elección del modelo base.

\begin{figure}[H]
\centering
\includegraphics[width=1\textwidth]{Resource/rq4-1.png}
\caption{Impacto de las estrategias de prompting en el rendimiento
         de LLMs para tareas de ciberseguridad, según valores
         representativos del corpus SLR. Fuente: Elaboración propia
         a partir de datos del corpus.}
\label{fig:impacto_prompting}
\end{figure}


\subsection{Interpretabilidad y XAI en análisis de logs}

LLM4Sec~\cite{karlsen2023llm4sec} establece el precedente de XAI aplicado
a logs de seguridad mediante TranSHAP (extensión de Kernel SHAP para
transformers). LogPrompt~\cite{liu2024logprompt} formaliza cinco niveles
de interpretabilidad (escala Likert 1--5) como criterios de evaluación
replicables.

KLAGE~\cite{belcastro2026klage} avanza en este frente mediante LIME para
explicabilidad post-hoc a nivel de nodo en Knowledge Graphs: los reportes
con contexto estructurado (GPT-adv) son preferidos en el 60\,\% de los
casos frente a los reportes sin contexto (GPT-base, preferido solo en el
7\,\%). Esta evidencia cuantitativa refuerza que la trazabilidad del
razonamiento no es un complemento cosmético sino un requisito funcional.

Palma et al.~\cite{palma2025llm} cuantifican por primera vez el drift
en la calidad explicativa de LLMs locales: la deriva temporal del FNR
es de solo 1.0\,\% a 3.5\,\% en 12 meses.

\subsection{Correlación multi-fuente y razonamiento causal}

Hmimou et al.~\cite{hmimou2025multiagent} demuestran que el razonamiento
multi-agente con correlación \textit{cross-domain} mejora la precisión: el
Cross-Context Recommendation System alcanza precisión de correlación de
amenazas del 87.0\,\% y\textit{ mean analyst trust score} de 4.6/5.0. MAD-CTI~\cite{shah2025madcti}
valida el patrón multi-agente para CTI: la generación de reportes de
amenazas con trazabilidad de razonamiento
es el componente más valorado por los analistas.

IVAM~\cite{ismail2025ivam} provee el diagrama de flujo end-to-end más
completo de un agente LLM para SOC, representado en la Tabla~\ref{tab:flujo_end_to_end}. 

\begin{table}[H]
\centering
\caption{Etapas del flujo end-to-end del agente.}
\label{tab:flujo_end_to_end}
\begin{tabular}{p{0.05\textwidth} p{0.35\textwidth} p{0.55\textwidth}}
\hline
\textbf{Etapa} & \textbf{Nodo} & \textbf{Función principal} \\
\hline
1 & \texttt{router\_node} & Recibe la entrada inicial y define la ruta de procesamiento del agente. \\

2 & \texttt{incident\_similarity\_node} & Recupera incidentes similares o casos históricos relevantes para contextualizar el análisis. \\

3 & \texttt{extracted\_problem\_node} & Extrae y estructura el problema principal, identificando hallazgos clave en los logs. \\

4 & \texttt{technical\_steps\_node} & Genera pasos técnicos de análisis, validación o remediación. \\

5 & \texttt{human\_interrupt\_node} & Incorpora una revisión humana para validar, ajustar o aprobar el flujo antes del cierre. \\

6 & \texttt{report\_node} & Construye el informe final
\hline
\end{tabular}
\end{table}


%%-------------------------------------------------------------
\section{Brechas identificadas y posicionamiento de la propuesta}
%%-------------------------------------------------------------

La revisión de los 52 trabajos del corpus permite identificar tres
brechas sistemáticas que fundamentan la propuesta de este seminario:

\textbf{Brecha 1 — Ausencia de razonamiento causal iterativo sobre logs
heterogéneos.} La totalidad de los trabajos de análisis de logs revisados
(RQ1, RQ2) opera como \textit{pipeline }lineal pasivo: ingesta logs, produce una
clasificación o resumen, y termina. Ninguno implementa un ciclo de
razonamiento agentivo donde el sistema formule una hipótesis, la consulte
activamente contra una base de conocimiento, la refine con nueva evidencia
y emita una conclusión iterativamente fundamentada. Esta brecha es señalada
explícitamente por Partovian et al.~\cite{partovian2024sms}, que reporta
la «incapacidad de identificar la causa raíz del problema» como el desafío
de capacidad más frecuente en la literatura, y por García Gómez et al.~\cite{garcia2026collaborative},
que identifican el razonamiento causal como la capacidad más ausente en
los \textit{frameworks }LAD actuales.

\textbf{Brecha 2 — Falta de integración con inteligencia de amenazas
estructurada.} Solo tres trabajos del corpus integran MITRE ATT\&CK en
su pipeline: LogPrécis~\cite{boffa2024logprecis} lo hace acotado a NER
sobre shell logs; IVAM~\cite{ismail2025ivam} lo usa para mapeo de TTPs
sin construir un RAG sobre la base completa; y LogRESP-Agent~\cite{lee2025logrespagent}
lo integra via TTPMapper pero sin poblar la base vectorial desde STIX.

\textbf{Brecha 3 — Ausencia de narrativas de incidente completas y
accionables.} Los trabajos más cercanos a la generación de texto explicativo
(e.g. LogPrompt, KLAGE, Hmimou et al., CyberRAG) producen explicaciones de
predicciones individuales, no narrativas que cubran la cadena completa de
eventos con
secuencia temporal, técnicas MITRE, niveles de confianza por hallazgo y
trazabilidad del razonamiento.

La Tabla~\ref{tab:brechas} sintetiza las tres brechas identificadas,
los trabajos del corpus que las evidencian y la respuesta que propone
este seminario.

\begin{table}[H]
\centering
\caption{Síntesis de brechas identificadas en el corpus SLR y
         posicionamiento de la propuesta.}
\label{tab:brechas}
\renewcommand{\arraystretch}{1.4}
\begin{tabularx}{\textwidth}{p{0.5cm} p{3.2cm} p{3.8cm} X}
\toprule
\textbf{\#} & \textbf{Brecha} & \textbf{Evidencia en corpus} & \textbf{Respuesta del seminario} \\
\midrule
1 &
  Ausencia de razonamiento causal iterativo &
Partovian et al.~\cite{partovian2024sms}; García Gómez et al.~\cite{garcia2026collaborative} &
Arquitectura \textit{Plan-and-Execute} con LangGraph: formulación de hipótesis, planificación del análisis, consulta a bases de conocimiento y refinamiento progresivo de la reconstrucción del incidente \\
2 &
  Falta de integración con CTI estructurada &
  Solo 8/52 trabajos integran MITRE ATT\&CK; ninguno puebla RAG desde STIX completo &
  ChromaDB poblado desde repositorio oficial MITRE ATT\&CK CTI (STIX~2.1) \\
3 &
  Ausencia de narrativas de incidente completas &
  LogPrompt, KLAGE, CyberRAG producen explicaciones individuales sin cadena causal &
  Informe estructurado: línea temporal, TTPs, confianza por hallazgo y trazabilidad \\
\bottomrule
\end{tabularx}
\end{table}

